\subsection{Single-Repository-Strategie}
\label{subsec:single-repository-strategie}

Die fünf Profile sind keine Template-Forks. Core, Feature-Implementierungen,
Profile und Generator liegen in einer Wartungsbasis, aus der eigenständige
Produktprojekte entstehen
\parencite{kleiveist2026profiles,kleiveist2026template}.

Das Master enthält auch optionale Implementierungen und aktiviert selbst
PostgreSQL. Ein generiertes Projekt erhält dagegen nur Core, aufgelöste
Feature-Pfade, Manifest, Frontend-Profil und passende \path{.env.example}. Es
bleibt danach vom unveränderten Master getrennt. Die zentrale Pflege reduziert
parallele Baselines; eine automatische Rücksynchronisierung generierter
Projekte besteht nicht.
Diese Trennung schützt die Wartungsbasis vor Produktänderungen. Umgekehrt
profitieren bereits erzeugte Produkte nicht ohne gezielte Übernahme von
späteren Korrekturen des Masters.
\beginnerinput{workspace/statement/400_projektprofile/480}
